
\section{QO-03: Beam splitters, HOM interference and NOON states}
\source{Week 2 Tutorial, Questions 1--2; Optics Assignment 1 Question 2;
2025 Exam Question 3; Week 3 Tutorial worked sheet.}
\subsection*{Step 1: select one convention and keep it}
Use the balanced transformation from the past exams,
written in terms of input creation operators,
\[
\boxed{\hat a^\dagger\longrightarrow
\frac{\hat c^\dagger+\hat d^\dagger}{\sqrt2},\qquad
\hat b^\dagger\longrightarrow
\frac{\hat c^\dagger-\hat d^\dagger}{\sqrt2}.}
\]
The four-port device conserves photon number because the
mode matrix is unitary. For a general real splitter with
\(t^2+r^2=1\), \(a^\dagger\to t c^\dagger+r d^\dagger\),
\[
\ket{n,0}\longrightarrow\sum_{k=0}^n
\sqrt{\binom nk}\,t^k r^{n-k}\ket{k,n-k}.
\]
Thus \(P(k,n-k)=\binom nk t^{2k}r^{2(n-k)}\),
the familiar binomial distribution.

\subsection*{Step 2: expand each two-photon input}
Use \(\ket{2,0}=(a^\dagger)^2\ket0/\sqrt2\).
The normalised output states are
\begin{align*}
\ket{2,0}_{ab}&\longrightarrow
\tfrac12\ket{2,0}_{cd}
+\tfrac1{\sqrt2}\ket{1,1}_{cd}
+\tfrac12\ket{0,2}_{cd},\\
\ket{0,2}_{ab}&\longrightarrow
\tfrac12\ket{2,0}_{cd}
-\tfrac1{\sqrt2}\ket{1,1}_{cd}
+\tfrac12\ket{0,2}_{cd},\\
\ket{1,1}_{ab}&\longrightarrow
\frac{\ket{2,0}_{cd}-\ket{0,2}_{cd}}{\sqrt2}.
\end{align*}
The last line is the Hong--Ou--Mandel result: indistinguishable
photons entering opposite ports never produce a coincidence.

\subsection*{Step 3: solve the 2025 NOON input}
For \(\ket{\Psi_{\rm in}}=(\ket{2,0}+\ket{0,2})/\sqrt2\),
the opposite-sign \(\ket{1,1}\) amplitudes cancel:
\[
\boxed{\ket{\Psi_{\rm out}}=
\frac{\ket{2,0}+\ket{0,2}}{\sqrt2},\quad
P_{20}=P_{02}=\tfrac12,\quad P_{11}=0.}
\]
The plus NOON state is invariant under this particular balanced
beam splitter. The \emph{minus} NOON input would instead
become \(\ket{1,1}\).

For the alternative 2025 input
\(\ket{\Phi_{\rm in}}=(\ket{1,1}+\ket{2,0})/\sqrt2\),
only the second component creates \(\ket{1,1}_{cd}\):
its amplitude is \((1/\sqrt2)(1/\sqrt2)=1/2\), so
\[
\boxed{P_{11}=\left|\tfrac12\right|^2=\tfrac14.}
\]
\checkmethod{Always convert powers of creation operators back
into \emph{normalised} Fock states; the \(\sqrt{2!}\) factors
change probabilities. Compare amplitudes before squaring them.}
